\section{本章小结}
\label{sec:fcover-summary}

本章开头问：能否预先造一张有限的函数表，使一类函数中的每个成员都在给定精度内有一行与之对应，这张表又需要多少行。用覆盖数表述，前一问是$N(\varepsilon,\mathcal F)$是否对每个$\varepsilon$有限，后一问是它随$\varepsilon$怎样增长。表中的行可以取在环境空间中，也可以要求取在类内，两种约定只差把精度减半。这个问题与为单个函数挑选近似不同：后者的预算可以随成员变化，前者要求一份预算服务整类。答案也依赖距离，同一族指示函数在积分距离下可以有限覆盖，在一致距离下不能。

只限制范数不足以得到有限覆盖。无限维空间的闭球中有彼此远离的序列，$C[0,1]$中越来越陡的帽形函数就是例子。在紧度量空间上，Arzelà--Ascoli判据给出了一致距离下的完整回答：连续函数类全有界、相对紧、一致有界且等度连续三者等价，类自身紧还需闭性。证明的机制是有限化，等度连续让有限个位置代表整个定义域，一致有界让每个位置只有有限种取值记录。紧定义域不能省略，$\mathbb R$上的平移族共用Lipschitz常数，却没有有限覆盖。

把这一机制做成具体构造，就得到计数。对$[0,1]^d$上幅值不超过$M$、Lipschitz常数为$L$的函数，网格误差由$L/m$控制，取值误差由量化间距控制，多线性插值把节点值连成连续的代表。直接计数给出度量熵上界$(\lceil2L/\varepsilon\rceil+1)^d\log\lceil2M/\varepsilon\rceil$；利用相邻节点的约束可以去掉主项中的对数因子，符号隆起的打包集又给出$(L/(6\varepsilon))^d\log2$的下界。维数出现在指数上，幅值只进入对数项。另一条路线是有限维表示：整类若能在同一个$n$维子空间中以误差$\varepsilon_0$统一逼近且近似有界，度量熵就不超过$n\log(1+2R/\varepsilon_1)$；按系数量化时还要核对系数界，允许从多个子空间中选择时要为候选个数付出$\log J$。

本章的结论都是确定性的：给定函数类、距离与精度，代表表是否存在、有多大。把它们用于由随机样本估计函数的问题，还需要另外的概率工具。下一章转向周期函数和全空间函数的频率表示，其中Fourier系数的衰减速度正是本章正弦级数类例子所用的那类条件；怎样构造满足统一逼近前提的有限表示，则在第~\ref{chap:function-approximation}章研究。
